\chapter{Conclusion}
\label{ch:conclusion}

This thesis addressed one limitation of the software project management platforms in wide use today: because a single provider stores and controls a project's history, nobody else can verify that the history has not been tampered with (\Cref{sec:motivation}). Starting from the architectural pattern that \textcite{adel2023} validated for decentralized data management in construction (\Cref{sec:adel}), the thesis adapted that pattern to software project management and implemented it in a working prototype, ChainBoard.

ChainBoard models accounts, projects and tasks as chaincode-governed state on a permissioned Hyperledger Fabric network. Its operations for creating projects, managing membership and roles, assigning tasks, changing status, commenting and archiving are all recorded in Fabric's immutable transaction history, not in a log controlled by a provider (\Cref{sec:chaincode-design}). The chaincode itself decides who may do what, from the roles stored in the project, so the permission rules do not depend on the backend or the user interface (\Cref{sec:permissions}). File attachments are stored on IPFS and anchored to the ledger by their content identifiers, which keeps binary data off the chain while still tying each recorded reference to the exact file content (\Cref{sec:ipfs-integration}). A backend API and a browser frontend offer these functions through a task board similar to Trello or Jira (\Cref{sec:backend,sec:frontend}). The audit trail, the main guarantee of the system, is visible and can be checked: every change shows who made it and when, and a verification step checks each record against the ledger and produces a downloadable report with a digest chain that can be verified again offline (\Cref{sec:audit-verification}).

The evaluation in \Cref{ch:results} showed that the design works as intended. Project and task operations behave as specified and the chaincode enforces the role rules; the audit trail keeps every past state of a task or project and can be verified; and file attachments are anchored by the CID recorded on-chain and can be retrieved from IPFS with it. The evaluation also reports what these guarantees cost and what they do not cover. Overall, the Fabric-plus-IPFS pattern, so far validated only in construction, carries over to software project management with the same architecture; what changes is the data model and the workflow offered to users.

\section{Answers to the Research Questions}
\label{sec:rq-answers}

\textbf{RQ1.} Yes. The Fabric-plus-IPFS pattern carries over to software project management with a different data model (users, projects, tasks, comments and attachments, with keys designed for a store that has no rich queries) and a familiar board workflow. The prototype meets all the requirements R1--R9 of \Cref{tab:requirements}.

\textbf{RQ2.} Yes, with one qualification. The chaincode checks the role of the acting person on every state-changing call, so the rules hold even if the backend or the interface is bypassed, and refused calls leave no trace on the ledger (\Cref{sec:results-permissions}). The qualification is that the chaincode learns \emph{who} is acting from the backend, because Fabric authenticates the organization and not the person. Closing this gap requires one Fabric identity per person.

\textbf{RQ3.} The system checks every transaction against the peer, computes a digest chain over the history, and lets the user download a report that a stand-alone script verifies offline. Together these let a person check a history without trusting the application that shows it (\Cref{sec:audit-verification,sec:results-audit}). The check shows that a history was not altered. It does not show that the person recorded as acting was really the one who acted.

\textbf{RQ4.} A write costs about the ordering service's batch timeout plus endorsement and commit, reads are answered by a peer and are much faster, and concurrent edits of the same key conflict (\Cref{sec:results-performance}). The limits are listed in \Cref{tab:threats} and in the next section: identity, confidentiality of reads, permanence of personal data, availability of attachments, and the scale of the evaluation.

\section{Limitations}
\label{sec:limitations}

The prototype was evaluated on a local Fabric network with two organizations, not on a geographically distributed or production-scale deployment. The performance figures in \Cref{ch:results} describe one machine under a synthetic load, not realistic use by many people, and they include no comparison with a conventional database. The most important limitation is the identity model. Fabric's gateway authenticates the organization and not each person, so the backend authenticates people with its own login sessions and tells the chaincode who is acting. The chaincode enforces the \emph{rules} but has to trust the backend about \emph{who} is acting, and a compromised backend could claim any identity (\Cref{sec:results-threats}). Reads are not restricted by role, so any signed-in person who knows a project identifier can read it. Data written to the ledger is permanent, which conflicts with the right to erasure of the GDPR~\cite{gdpr} for personal data such as names and comments. Attachments are anchored by content identifier, but nothing in the prototype guarantees that the files stay available. The data model is deliberately small. It supports a Kanban-style board but not the Scrum framework (there is no backlog ordering, no sprints, no estimation and no burndown chart; see \Cref{tab:agile}), and it does not yet cover sub-tasks, dependencies between tasks, notifications or access control per field. Finally, the determinism constraint of \Cref{sec:determinism} and Fabric's optimistic concurrency show that systems of this kind have development costs that centralized platforms do not, because the code must stay correct when several peers execute it.

\section{Future Work}
\label{sec:future-work}

The limitations above suggest several directions. The most useful is to give every person their own Fabric identity, with a certificate issued by the Fabric certificate authority, so that the chaincode reads the caller from the transaction's client identity instead of trusting the backend, as the Fabric-based tool of \textcite{batista2021} already does, and to require endorsement by both organizations for critical actions such as archiving a project. Attribute-based policies of the kind studied by \textcite{gordijn2022combining} could then express the project roles directly in the certificates. Restricting reads by role, or placing sensitive fields in Fabric private data collections and keeping personal data off-chain with only hashes on the ledger, would address confidentiality and erasure. Evaluating the system on a geographically distributed network with a realistic number of organizations and peers, and comparing it with a conventional database, would show its throughput and latency under conditions closer to production. Extending the data model with sprints, a backlog, story points, burndown charts, sub-tasks, dependencies and notifications would close the functional gap with mature centralized platforms while keeping the same audit guarantees, and pinning attachments, with a check of the content identifier on download, would make file storage as dependable as the ledger. Finally, the architectural pattern is not specific to software projects. Applying the same chaincode-plus-IPFS approach to other collaborative domains that need a verifiable history would show how far it generalizes beyond construction and software project management, the two domains where it has now been demonstrated.
